\documentclass{situationist-preprint}

\title{Bundle Theory Refuted}
\msid{MET-29-a}
\msfamily{MET-29}
\msdiscipline{Metaphysics}
\mstype{Principal manuscript}
\msstatus{Revised 7 October 2026}
\mskeywords{bundle theory; identity of indiscernibles; weak discernibility; substratum; individuation}

\begin{document}
\maketitle

\begin{abstract}
The bundle theory holds that an object is nothing over and above a bundle of qualities. It entails a version of the identity of indiscernibles: objects with the same qualities are the same bundle, and so the same object. Given the results of MET-28, bundle theories that rely on absolute or relative discernibility are refuted by Black's spheres. A bundle theory that allows relations among bundles to individuate them, and so relies only on weak discernibility, survives.
\end{abstract}

\paragraph{Note on the revision.} The version of 6 October 2026 claimed that every bundle theory is refuted. Following the revision of MET-28-a, the claim is restricted to bundle theories based on absolute or relative discernibility.

\section{The argument}

The bundle theory, in the form defended by Hume and by \citet{Russell1948}, identifies an object with the bundle of its qualities. If two objects had exactly the same qualities, they would be the same bundle, and the theory would count them as one. The theory is therefore committed to a principle of the identity of indiscernibles, and its fate depends on which version of that principle is true \citep{OLearyHawthorne1995}.

\begin{proposition}\label{p:main}
A bundle theory on which bundles are individuated by their monadic qualities, or by their qualitative relations taken in order, is false.
\end{proposition}

\begin{proof}
Such a theory entails absolute or relative PII. By MET-28-a, Theorems~3.1 and 3.2, both principles fail in Black's world, which contains two distinct spheres with the same monadic qualities and the same qualitative relations in either order.
\end{proof}

\begin{proposition}\label{p:weak}
A bundle theory on which bundles may be individuated by irreflexive qualitative relations among them is not refuted by Black's world.
\end{proposition}

\begin{proof}
By MET-28-a, Proposition~4.1, the spheres are weakly discernible. A bundle theory that admits weak discernibility counts them as two.
\end{proof}

Whether the surviving theory is still a bundle theory is open to question. If what makes the two spheres two is that they stand in a relation of distance to each other, the relation does individuating work that the qualities alone cannot do, and the theory has moved toward structuralism. We note the point without pursuing it.

\begin{thebibliography}{9}\small
\bibitem[O'Leary-Hawthorne(1995)]{OLearyHawthorne1995} O'Leary-Hawthorne, J. (1995). The bundle theory of substance and the identity of indiscernibles. \emph{Analysis}, 55(3), 191--196.
\bibitem[Russell(1948)]{Russell1948} Russell, B. (1948). \emph{Human Knowledge: Its Scope and Limits}. Allen \& Unwin.
\end{thebibliography}

\end{document}
